This section discusses potential mechanisms behind the observed window length and shape effects.
The discussion is hypothesis driven.
The goal is to guide future experiments that test these mechanisms directly.

Length tuples control which time scales are exposed to the model.
Click behavior can change over minutes, hours, and days.
Short windows can capture immediate recency effects and repeated exposure.
Day and week scale windows can capture periodicity and slower drift.
If a length tuple omits a relevant scale, a different window shape cannot recover it.
In that case, shape changes mainly redistribute the same historical mass.

Trailing windows represent recent history and its decay in a direct way.
They also provide a simple bias variance tradeoff.
As $L$ increases, a trailing window reduces variance of rate estimates but can increase bias under drift.
A small set of trailing windows creates a multi resolution summary.
A tree model can combine these summaries in a nonlinear way.
This can be sufficient when the main signal is recency and repetition.

Gap windows exclude the most recent hour from the history window.
This can remove the most informative part of the signal when behavior changes quickly.
It can also remove short term dependence for repeated identifiers.
Gap windows can be useful when very recent data is contaminated by a leakage channel.
In this study, the no lookahead constraint already excludes same hour leakage.
Therefore, the gap mainly acts as an unnecessary information restriction.

Bucketized windows partition history into disjoint intervals.
This increases feature dimensionality and can increase estimation variance.
Each bucket contains fewer events than an overlapping trailing window at a similar horizon.
Bucketization also removes cumulative smoothing from nested trailing windows.
If the signal is approximately monotone in recency, trailing windows represent it efficiently.
Bucketization may help when distinct time bands have distinct effects.

Fixed length time windows implicitly assume a stable event rate.
For many entity keys, event rates vary widely across entities and over time.
An event count window adapts by fixing the number of observations used to estimate a rate.
This can stabilize estimates for sparse entities.
It can also reduce the impact of bursty traffic for frequent entities.
It may also be more robust under hour resolution timestamps with irregular sampling.

Calendar aligned windows aim to represent full previous days or weeks.
They can help when boundary effects exist at calendar transitions.
However, in hour resolution data, calendar alignment can be redundant with explicit time features.
Examples include hour of day and day of week.
Calendar blocks can also mix heterogeneous periods under drift.
These factors can yield small effects with signs that vary across length tuples.
